% !TEX program = xelatex
% 肝癌空间免疫抑制微环境分析项目报告（Nature 风格重构版）
% 使用 ctexart 文档类，XeLaTeX 编译

\documentclass[12pt, a4paper]{ctexart}

% ============================================================
% 宏包加载
% ============================================================
\usepackage{geometry}
\geometry{left=2.5cm, right=2.5cm, top=2.5cm, bottom=2.5cm}

\usepackage{graphicx}
\usepackage{booktabs}
\usepackage{longtable}
\usepackage{amsmath}
\usepackage{amssymb}
\usepackage{hyperref}
\usepackage{xcolor}
\usepackage{listings}
\usepackage{float}
\usepackage{caption}
\usepackage{subcaption}
\usepackage{enumitem}
\usepackage{setspace}
\usepackage{fancyhdr}
\usepackage{natbib}
\usepackage{array}
\usepackage{multirow}
\usepackage{tabularx}
\usepackage{appendix}
\usepackage{microtype}

% ============================================================
% 页面与格式设置
% ============================================================
\onehalfspacing
\pagestyle{fancy}
\fancyhf{}
\fancyhead[L]{\small 基于单细胞与空间转录组的肝癌免疫抑制微环境分析}
\fancyhead[R]{\small \thepage}
\renewcommand{\headrulewidth}{0.4pt}

\lstset{
    basicstyle=\small\ttfamily,
    keywordstyle=\color{blue},
    commentstyle=\color{gray},
    stringstyle=\color{red},
    breaklines=true,
    frame=single,
    numbers=left,
    numberstyle=\tiny\color{gray},
    language=Python,
    showstringspaces=false,
    xleftmargin=1em,
    framexleftmargin=1em,
    extendedchars=false,
    keepspaces=true
}

\hypersetup{
    colorlinks=true,
    linkcolor=blue,
    citecolor=blue,
    urlcolor=blue
}

% ============================================================
% 标题信息
% ============================================================
\title{
    \vspace{-1cm}
    \textbf{基于单细胞与空间转录组整合的肝癌\\
    免疫抑制微环境空间解析与预后验证}\\[0.5em]
    \large 项目技术报告
}
\author{}
\date{\today}

% ============================================================
% 正文开始
% ============================================================
\begin{document}

\maketitle
\thispagestyle{empty}

\begin{abstract}
肝细胞癌（Hepatocellular Carcinoma，HCC）是全球癌症相关死亡的第三大原因，
免疫检查点抑制剂的临床有效率仅约15\%--20\%，
核心瓶颈在于肿瘤微环境中调节性 T 细胞（Treg）、
肿瘤相关巨噬细胞（TAM）与肿瘤相关成纤维细胞（CAF）所共同构建的多层次免疫抑制网络~\citep{llovet2021hepatocellular,finn2020atezolizumab}。
本项目整合单细胞转录组（scRNA-seq）与 10x Genomics Visium 空间转录组数据，
以 Cell2location 贝叶斯反卷积算法~\citep{kleshchevnikov2022cell2location}
将细胞类型注释精确映射至空间切片的每个 spot，
并采用基于 Schürch 等人"细胞邻域（Cellular Neighborhoods）"
框架~\citep{schurch2020coordinated}的邻域组成聚类策略，
在不依赖先验阈值的前提下识别免疫抑制空间生态位。
通过配体-受体通讯定量分析~\citep{efremova2020cellphonedb}
揭示 Treg--TAM--CAF 轴核心信号通路（CXCL12--CXCR4、CCL22--CCR4、TGFB1--TGFBR1、PD-1--PD-L1），
利用两张独立 Visium 切片（HCC4R 与 CHC20）验证跨样本一致性，
并以 ssGSEA 方法将提取的空间特征基因签名投影至 TCGA-LIHC 队列
（约370例患者）~\citep{tomczak2015cancer}进行独立预后验证。
本项目建立了从空间精细解析到机制性通讯分析再到临床队列验证的完整转化分析框架。

\noindent\textbf{关键词：}
肝细胞癌；空间转录组；单细胞转录组；免疫抑制微环境；Cell2location；
Cellular Neighborhoods；配体-受体通讯；ssGSEA；生存分析
\end{abstract}

\newpage
\tableofcontents
\newpage

% ============================================================
%  第一部分：研究背景与意义
% ============================================================

\section{引言}

\subsection{疾病背景与临床挑战}

肝细胞癌（HCC）是全球第六大恶性肿瘤，每年新发病例超过90万，死亡病例约83万，
其中中国约占全球新发病例的45\%~\citep{llovet2021hepatocellular,chen2021liver}。
由于大多数患者确诊时已属晚期，HCC 的五年生存率不足20\%。
索拉非尼曾是晚期 HCC 的标准一线方案，但相比对照组仅将中位总生存期延长约3个月~\citep{llovet2008sorafenib}。
2020年 IMbrave150 试验证明，阿替利珠单抗联合贝伐珠单抗可将客观缓解率提升至27.3\%，
并首次在总生存期上优于索拉非尼~\citep{finn2020atezolizumab}；
然而仍有约70\%的患者对该联合方案无应答，揭示肿瘤微环境中存在强大的多层次免疫抑制机制。

HCC 微环境中的免疫抑制由多种细胞协同构建。调节性 T 细胞（Treg）以 CD$4^+$CD$25^+$ FOXP$3^+$ 为经典表型，通过分泌 IL-10、TGF-$\beta$ 和 IL-35，以及表达 CTLA-4、TIGIT、LAG-3 等检查点分子，
直接抑制效应性 T 细胞与 NK 细胞的抗肿瘤活性~\citep{sakaguchi2020regulatory}；
单细胞转录组研究表明 Treg 在肝癌组织中显著富集，
Treg 与 CD$8^+$ T 细胞的比值是独立预后因子~\citep{zheng2017landscape}。
以 M2 极化表型（CD163$^+$MRC1$^+$）为主的肿瘤相关巨噬细胞（TAM）通过 CCL22--CCR4 和 CXCL12--CXCR4 轴
主动招募 Treg 进入肿瘤区域，同时分泌 TGF-$\beta$ 并高表达 PD-L1 以抑制效应性 T 细胞~\citep{biswas2012macrophage,noy2014tumor}；
TAM 与 Treg 之间的双向正反馈循环共同维持了稳定的局部免疫抑制生态。
以 FAP、ACTA2 和 POSTN 为标志的肿瘤相关成纤维细胞（CAF）则通过分泌 CXCL12 招募免疫抑制细胞，
并通过重塑细胞外基质（ECM）形成物理屏障阻碍效应性 T 细胞浸润~\citep{kalluri2016biology,mariathasan2018tgfb}。
这一 Treg--TAM--CAF 三角轴的协同作用是 HCC 免疫治疗耐药的核心机制，
也是本项目重点解析的生物学对象。

\subsection{空间转录组学与整合分析范式}

上述研究需要理解免疫抑制的空间组织结构。
传统 bulk RNA-seq 将组织均质化后测序，丢失了细胞间的空间位置信息；
scRNA-seq 虽能解析单细胞转录特征，但在组织解离过程中不可避免地破坏了空间关系~\citep{luecken2019current}。
10x Genomics Visium 空间转录组平台在保留组织原位结构的前提下，
对每个空间位点（spot，直径约55\,$\mu$m）进行全转录组测序，
使研究者得以直接观测细胞类型的空间分布模式~\citep{williams2022introduction}。
由于单个 spot 覆盖1--10个细胞的混合信号，
需结合 scRNA-seq 参考数据通过计算反卷积对其细胞组成进行定量估计。

在当前主流反卷积工具中，Cell2location 基于负二项贝叶斯层次模型，
在全基因组范围内同时估计多种细胞类型的绝对丰度，并提供后验不确定性量化~\citep{kleshchevnikov2022cell2location}；
多项独立基准测试一致表明，它在稀有细胞类型（如 Treg，比例通常 $<$5\%）检测上
优于 RCTD~\citep{cable2021robust}、CARD~\citep{ma2022integrative} 和 Tangram~\citep{biancalani2021deep}~\citep{li2022benchmarking,sang2023spotless,yan2023benchmarking}。
Schürch 等人在结直肠癌中发展的"细胞邻域"概念，
通过对局部细胞组成进行无监督聚类揭示不同 neighborhood 对抗肿瘤免疫的差异化贡献，
为空间分辨率下的肿瘤微环境分型提供了方法学范式~\citep{schurch2020coordinated}。
在此基础上，将空间转录组局部特征基因签名通过 ssGSEA~\citep{barbie2009systematic}
投影至大规模临床队列，结合多变量 Cox 回归~\citep{cox1972regression}
进行独立预后验证，是当前实现"空间发现--临床转化"的经典策略~\citep{hanzelmann2013gsva}。

现有 HCC 研究仍普遍缺少三个维度的整合：
在保留空间信息的前提下，将 Treg--TAM--CAF 轴的共定位从描述性层面提升至配体-受体通讯机制层面；
采用数据驱动的方法识别免疫抑制空间生态位；
以及将空间转录组发现系统性地在独立大规模临床队列中进行预后验证。
本项目针对上述不足，建立了一套可推广的完整转化分析框架。

\subsection{研究目标}

本项目旨在回答四个相互关联的核心科学问题：
（1）肝癌 Visium 切片中，免疫抑制细胞信号（Treg、TAM、CAF）是否存在系统性空间共定位，
且富集于特定的肿瘤-基质交界区域？
（2）该区域内 Treg--TAM--CAF 轴是否通过特定配体-受体对（CXCL12--CXCR4、CCL22--CCR4 等）
实现协同免疫抑制？
（3）所识别的免疫抑制空间生态位特征是否在多张独立切片间具有跨样本一致性？
（4）基于空间数据提炼的特征基因签名能否在 TCGA-LIHC 大队列中独立预测患者预后？

% ============================================================
%  第二部分：方法
% ============================================================

\section{方法}
\subsection{数据来源与质量控制}

本研究整合了单细胞转录组（scRNA-seq）、空间转录组（ST）及大规模临床大规模转录组等多模态数据集，以构建“外周发现--空间定位--临床验证”的转化医学分析框架。

首先，单细胞参考数据集下载自 GEO 数据库（登录号：GSE149614），涵盖肝细胞（Hepatocyte）、T/NK 细胞、B 细胞、髓系细胞（Myeloid Cell）和成纤维细胞（Fibroblast）等核心细胞类型。参照当前单细胞组学分析的最佳实践~\citep{luecken2019current}，本研究严格过滤了总表达量低于 200 的低质量细胞以及在少于 10 个细胞中表达的基因。鉴于后续所采用的 Cell2location 空间反卷积模型基于贝叶斯负二项回归框架，其内部已实现了对测序深度和技术批次效应的概率建模，因此在单细胞预处理阶段不执行常规的文库大小归一化或 $\log(1+p)$ 变换，直接保留原始 Count 矩阵作为模型的输入。

本研究所使用的空间转录组测序数据均由 10x Genomics Visium 高通量平台生成，采用双中心、跨临床情境的严谨设计构建空间微环境特征的发现与验证闭环。其中，作为主发现集的切片 HCC4R源自前瞻性临床试验队列（GSE238264），该样本来自接受卡博替尼联合纳武利尤单抗靶向免疫新辅助治疗后的HCC切除标本，具有丰富的HCC-CAF空间交互背景；独立外部验证集切片 CHC20源自另一项聚焦于非酒精性脂肪性肝炎相关肝癌的独立空间组学队列，该队列核心关注肿瘤纤维化病灶中骨髓细胞介导的免疫抑制机制。

参照空间转录组分析的最佳实践与质控标准~\citep{williams2022introduction,palla2022squidpy}，对上述两张 ST 切片进行统一过滤：仅保留总 UMI 计数在 5000--35000 之间、且线粒体基因比例低于 20\% 的高质量 Spots，同时剔除在少于 10 个 Spots 中表达的低频基因。

最后，为了将空间组学发现系统性地推向独立大规模临床队列进行生存预后验证，本研究引入了 TCGA-LIHC 队列（包含约 370 例 HCC 患者）。通过 TCGAbiolinks R 包~\citep{colaprico2016tcgabiolinks} 从 GDC 数据库下载标准化的 HTSeq-FPKM 格式表达矩阵以及完整的临床随访与生存信息。

\subsection{Treg 亚群鉴定与 Cell2location 反卷积}

从 scRNA-seq 数据中提取 T/NK 细胞亚群，
以 Leiden 算法（分辨率 res=3.0）进行高分辨率亚聚类~\citep{traag2019leiden}，
通过气泡图可视化 CD3D、CD4、FOXP3 和 IL2RA 在各子簇中的表达分布。
确认特定簇呈现 CD4$^+$CD25$^+$FOXP3$^+$ 的 Treg 典型特征后~\citep{sakaguchi2020regulatory}，
将该簇重注释为 Treg 并生成最终细胞类型标签。

Cell2location 分两阶段运行~\citep{kleshchevnikov2022cell2location}。
第一阶段，以 RegressionModel 在 scRNA-seq 数据（原始 count）上训练细胞类型参考签名矩阵，
最大训练轮次250 epochs，batch size 1024，early stopping patience=30；
第二阶段，以参考签名矩阵对每张 Visium 切片独立运行空间建模，
关键先验参数 N\_cells\_per\_location=30（基于 Visium spot 直径约55\,$\mu$m 及典型肝脏组织细胞密度估计~\citep{kleshchevnikov2022cell2location}），
detection\_alpha=20（官方推荐默认值），最大训练1000 epochs。
输出后验均值丰度矩阵（\texttt{means\_cell\_abundance\_w\_sf}）作为各 spot 细胞类型绝对丰度的估计。

\subsection{免疫抑制空间生态位的识别}

以 Cell2location 输出的绝对丰度归一化为比例矩阵后，
利用 $k$ 近邻算法（$k$-NN）在组织几何坐标上构建空间物理邻接图。
对每个 spot $i$ 计算其在物理空间中 $k$ 个最近邻 spot 的细胞类型比例均值，从而抽象出邻域组成向量 $\mathbf{n}_i$：
\begin{equation*}
\mathbf{n}_i = \frac{1}{|\mathcal{N}_i|} \sum_{j \in \mathcal{N}_i} \mathbf{p}_j
\label{eq:neighborhood_vector}
\end{equation*}
本研究沿用 Schürch 等人~\citep{schurch2020coordinated} 在结直肠癌研究中的经典参数选择 $k=15$，并通过 $k \in \{10, 15, 20\}$ 的敏感性分析验证空间拓扑结构的参数稳健性。

随后，以该邻域组成向量 $\mathbf{n}_i$ 为高维表型特征，在高维特征空间中重新构建 $k$-NN 图。最后，在此特征邻接图上运行 Leiden 社区检测算法（分辨率=0.5）以自适应识别空间生态位（Niche）类型。
与 KMeans 相比，Leiden 无需预先指定聚类数、且理论上保证了聚类连通性~\citep{traag2019leiden}。

对每个 Leiden 聚类，计算 Treg 比例、Myeloid 比例、
Fibroblast 比例和免疫抑制基因模块评分四个功能维度的均值，
按各维度降序排名后求和得到综合排名 $R_c$：
\begin{equation}
R_c = \text{rank}(\bar{p}_{\text{Treg},c}) + \text{rank}(\bar{p}_{\text{Myeloid},c})
+ \text{rank}(\bar{p}_{\text{Fibroblast},c}) + \text{rank}(\bar{s}_{\text{immune},c})
\label{eq:combined_rank}
\end{equation}
$R_c$ 最小的 cluster 注释为免疫抑制生态位（immunosuppressive niche），
该多维综合评分确保所识别的 niche 同时满足 Treg--TAM--CAF 三角轴的协同富集特征，
而非单一细胞类型的孤立高表达。

免疫抑制基因模块评分由14个文献支持的标志基因（FOXP3、IL2RA、CTLA4、TIGIT、LAG3、PDCD1、HAVCR2、TGFB1、IL10、CXCL12、CCL22、TNFRSF18、TNFRSF4、IKZF2）~\citep{sakaguchi2020regulatory,sharma2023immunosuppressive}
的 CP10K+log1p 表达均值定义；
辅助评分体系（Treg-like score、immune-stroma score 和 niche score）通过对相应组成维度进行 Z-score 标准化后加和计算（详见附录）。
{\color{red}敏感性分析采用 Spearman 秩相关系数（$\rho$）和连续型 Weighted Jaccard 相似度
评估 $k \in \{10, 15, 20\}$ 下评分排序的稳定性，
以 $\rho > 0.90$ 且 weighted Jaccard $>$ 0.80 作为参数稳健性的通过标准~\citep{spearman1904proof}。}

\subsection{特征基因签名的三层筛选}
{\color{red}写得更详细一点}
Visium spot 同时覆盖多种细胞类型（细胞稀释效应），
使 FOXP3 等 Treg 标志基因在大多数 spot 中的原始表达量为零，
导致均值 log2FC 被严重低估。
为此，本研究以三层筛选策略提取免疫抑制生态位特征基因签名。

第一层（Layer 1）：以空间网格聚合（block\_mwu 模式）将连续 spot 合并为 block 均值以降低空间自相关~\citep{palla2022squidpy}，
对 {\color{red}niche\_high（$\geq$85百分位）}vs niche\_low 组执行 Wilcoxon 检验+BH-FDR 校正~\citep{mann1947test}；
保留 FDR $<$ 0.1 且 AUC $>$ 0.60~\citep{hanley1982meaning}
且（log2FC $>$ 0.3 OR $\Delta$frac $>$ 0.10）的基因，
同时剔除涵盖管家基因、核糖体蛋白、线粒体基因等350余个非特异性基因的黑名单；
以 $\text{signature\_rank\_score} = 0.35 \cdot \text{minmax}(\log_2\text{FC}) + 0.35 \cdot \text{minmax}(\Delta\text{frac}) + 0.30 \cdot \text{minmax}(\text{AUC})$
对候选基因进行加权综合排序。
第二层（Layer 2）：针对 Treg 标志物（FOXP3、IL2RA、CTLA4 等）、
TAM 特征基因（CD163、MRC1、TGFB1 等）和 CAF 激活基因（FAP、ACTA2、COL1A1 等）
三组先验功能基因集分别执行 AUC 检验（阈值 AUC $>$ 0.60 且 $\Delta$frac $>$ 0.05），
强制纳入符合条件的先验基因以防止 Visium 稀疏性导致的生物学关键基因漏选。
{\color{red}第三层（Layer 3）：通过 Gini 不均等指数~\citep{yitzhaki1979relative}（阈值 Gini $>$ 0.3 且 log2FC $>$ 0）
捕获局灶性高表达的稀有免疫基因，
弥补前两层均值比较对高度稀疏基因的系统性低估。}
{\color{red}删除第三层，因为输出的gini指数表格是空的}

\subsection{配体-受体通讯分析}
{\color{red}公式写得更规范一点}

参考 CellPhoneDB~\citep{efremova2020cellphonedb} 和 COMMOT~\citep{cang2023commot} 数据库，
选取12对在 Treg--TAM--CAF 轴中具有明确生物学意义的配体-受体对
（CXCL12--CXCR4、TGFB1--TGFBR1、PD-1--PD-L1、CCL22--CCR4、
TIGIT--NECTIN2、LAG3--HLA-DRA、IL10--IL10RA、SPP1--CD44、
MIF--CD74、VEGFA--KDR、Galectin9--TIM-3、CCL2--CCR2；
完整列表见附录表~\ref{tab:lr_pairs}）。
对每对 $(l, r)$，以 CP10K+log1p 标准化表达量的乘积作为通讯强度代理指标，
$\text{LR\_score} = E_l \cdot E_r$，
该乘积形式与 CellPhoneDB 实现一致，自然地满足"配体和受体须同时存在"的生物学约束~\citep{efremova2020cellphonedb}。
分别计算 niche\_high 与 niche\_low 组的均值 LR\_score，
以 $\text{log2FC} = \log_2[(\bar{\text{LR}}_{\text{high}} + \varepsilon)/(\bar{\text{LR}}_{\text{low}} + \varepsilon)]$（$\varepsilon = 10^{-6}$）
量化生态位内的通讯富集倍数。
对主基因在数据集中缺失的情况，自动尝试同家族替代基因以最大化可分析通讯对的数量。

\subsection{跨切片一致性验证与 TCGA 生存分析}

以 HCC4R 提取的签名基因集为基准，将免疫模块基因（Treg/TAM 相关，共21个）
和基质/ECM 模块基因（CAF/ECM 相关，共21个）分别在 CHC20 切片上计算模块评分，
以 Z-score 标准化后加和得到组合 niche 评分（immune\_stromal\_niche\_score），
通过空间投影图和 Mann-Whitney U 检验~\citep{mann1947test}
量化评分分布的跨样本一致性。

生存分析使用 GSVA R 包~\citep{hanzelmann2013gsva} 的 ssGSEA 方法~\citep{barbie2009systematic}，
以 $\log_2(\text{FPKM}+1)$ 矩阵为输入，高斯核（kcdf=Gaussian）、绝对排名（abs.ranking=TRUE）
对 TCGA-LIHC 队列中每位患者计算标量富集评分。
ssGSEA 为每位患者独立生成一个反映生态位活性的连续评分，
更适合纳入预后回归模型，而非 GSVA 的跨样本相对活性估计~\citep{barbie2009systematic}。
构建多变量 Cox 比例风险回归模型~\citep{cox1972regression}，
以 ssGSEA 评分、年龄和肿瘤分期为协变量：
\begin{equation*}
h(t \mid \mathbf{x}) = h_0(t) \cdot \exp(\beta_1 s_{\text{ssGSEA}} + \beta_2 \text{age} + \beta_3 \text{stage})
\label{eq:cox}
\end{equation*}
控制年龄和分期（HCC 已知强效独立预后因子~\citep{llovet2021hepatocellular}）
后，若 ssGSEA 评分的 $\beta_1$ 仍显著，则证明免疫抑制生态位签名捕获了独立于临床分期以外的预后信息。
以 ssGSEA 评分中位数为阈值，绘制 Kaplan-Meier 生存曲线~\citep{kaplan1958nonparametric}，
log-rank 检验~\citep{mantel1966evaluation} 评估两组生存差异。

% ============================================================
%  第三部分：结果
% ============================================================

\section{结果}

\subsection{scRNA-seq 参考图谱构建与 Treg 亚群鉴定}

利用 GSE149614 单细胞转录组数据构建细胞类型参考图谱。
细胞类型 Marker 基因热图（Figure 1C）以 Z-score 形式展示各细胞类型的平均表达特征，
肝细胞/恶性细胞（ALB、APOA2、APOA1）、T/NK 细胞（CD3D、GNLY、GZMB）、
髓系细胞（LYZ、AIF1、HLA-DRA）、HSC/Fibroblast（COL1A1、ACTA2、RGS5）、
B 细胞（MS4A1、CD79A）和内皮细胞（PECAM1、CDH5）均呈现清晰的对角线高亮模式，
验证了细胞类型注释的可靠性~\citep{ma2021comprehensive}。
值得注意的是，Treg 标志基因（FOXP3、IL2RA）的 Z-score 信号相对较低，
这是稀有细胞类型（比例 $<$2\%）在均值热图中的固有局限，
而非注释失误~\citep{luecken2019current}。

针对 T/NK 亚群，以 res=3.0 进行高分辨率 Leiden 聚类并绘制横版气泡图（Figure 1D）。
特定簇（簇9）呈现 FOXP3 气泡最大且颜色最深、CD4 共表达、IL2RA 显著高于其他簇的组合模式，
符合 CD4$^+$CD25$^+$FOXP3$^+$ 经典 Treg 表型~\citep{sakaguchi2020regulatory}；
其余 T/NK 亚簇 FOXP3 近乎缺失，证明了 FOXP3 表达的亚群特异性。
将簇9重注释为 Treg，生成最终细胞类型标签，该结果以高置信度纳入 Cell2location 反卷积的参考签名中。

\subsection{Cell2location 反卷积揭示 HCC4R 的免疫细胞空间分布}

Cell2location 反卷积将单细胞水平细胞类型注释精确映射至 HCC4R 切片每个 spot。
H\&E 染色图（Figure 2A）清晰呈现了肿瘤实质区（核大深染、Hepatocyte 密集）
与基质/间质区（纤维化、炎症浸润）的形态学边界，为空间转录组结果提供了组织学参照。

细胞类型空间分布图（Figure 2B）揭示了与预期高度一致的空间格局：
Hepatocyte 高丰度区域集中于肿瘤实质，与 H\&E 图中核密集区域的空间吻合验证了反卷积的准确性；
Treg 呈局灶性富集分布（右偏分布），高丰度热点主要定位于肿瘤实质与间质的交界区域，
与 Treg 优先聚集于"肿瘤-免疫界面"的生物学预期一致~\citep{zheng2017landscape}；
Myeloid 细胞与 Treg 的空间分布具有明显共定位趋势，
呼应了 TAM 通过 CCL22--CCR4 轴主动招募 Treg 的信号机制~\citep{noy2014tumor}；
Fibroblast 则主要分布于围绕肿瘤实质的基质层，
反映 CAF 构建物理屏障的功能定位~\citep{kalluri2016biology}。

细胞类型间的 Pearson 相关分析（Figure 2C）进一步量化了这种空间关系：
Treg 与 Myeloid 呈显著正相关，Fibroblast 与 Myeloid 同样正相关，
而三者均与 Hepatocyte 呈负相关——免疫基质细胞主要富集于肿瘤-基质过渡带而非肿瘤核心，
与经典"免疫排斥型"HCC 的空间结构描述一致~\citep{llovet2021hepatocellular}。

\subsection{Cellular Neighborhoods 聚类识别免疫抑制空间生态位}

以邻域组成向量（$k=15$，公式~\ref{eq:neighborhood_vector}）为特征运行 Leiden 聚类（分辨率=0.5），
在 HCC4R 切片中识别出多个（5--12个）具有内部一致性的空间域（Figure 3A）。
聚类结果呈现连续性空间分区而非随机散点，
直接验证了无监督聚类能够有效捕获组织的空间微结构~\citep{schurch2020coordinated}。
肿瘤核心区以 Hepatocyte 为主导，基质区以 Myeloid 和 Fibroblast 混合为特征，
而 Treg、Myeloid 与 Fibroblast 三者共同高度富集的 cluster 在综合排名（公式~\ref{eq:combined_rank}）
中得分最低，被注释为 immunosuppressive niche，
其主要分布于肿瘤实质与基质的交界地带，与预期的免疫抑制热点位置高度吻合。

各 domain 细胞组成堆积条形图（Figure 3B）清晰展示了 immunosuppressive niche cluster 的组成特征：
Treg 比例在所有 cluster 中最高，Myeloid 处于高位，Fibroblast 高于肿瘤实质 cluster，
而 Hepatocyte 处于中等水平——这种多细胞类型协同富集充分体现了 Treg--TAM--CAF 三角免疫抑制轴的空间协同性，
有别于单纯的淋巴组织样区域或纤维化区域~\citep{schurch2020coordinated,sakaguchi2020regulatory}。

小提琴图（Figure 3C）定量比较了 immunosuppressive niche 与其余 cluster 在
immune-stroma score、immunosuppressive niche score 和 Treg-like score 上的分布差异。
三个评分均在 immunosuppressive niche cluster 中显著高于其他区域（Mann-Whitney U 检验，$p < 0.05$），
immunosuppressive niche score 的两组分布差异最为突出：
非 niche cluster 集中于低分区间，而 niche cluster 呈明显右尾偏移。
这一结果表明，Leiden 聚类以无监督的细胞组成数据所识别的 niche 与功能评分系统的高度一致性，
为生态位注释的可靠性提供了独立验证~\citep{mann1947test}。

\subsection{niche 特征基因签名的提取与功能注释}

niche\_high（$\geq$85百分位）与 niche\_low 组的差异基因火山图（Figure 4A）揭示了
免疫抑制生态位的分子特征。右上象限中显著上调的基因涵盖：
免疫抑制核心因子（TGFB1、IL10、CXCL12、CCL22）、
CAF/ECM 重塑成分（FAP、ACTA2、COL1A1、SPP1）
以及多种免疫检查点相关分子。
FOXP3 和 IL2RA 虽在 log2FC 上值偏低，但 $-\log_{10}(\text{FDR})$ 有一定上移，
这与 Visium 稀释效应导致 Treg 标志基因均值 FC 被低估的预期完全吻合，
Layer 3 Gini Index 筛选对此进行了专项补偿~\citep{luecken2019current}。
与之相对，左侧显著下调基因主要为 Hepatocyte 特征基因（ALB、APOA2 等），
表明 niche\_high 区域肿瘤实质成分相对减少，免疫基质细胞占主导。

niche 签名基因在各空间区域的表达热图（Figure 4B）进一步揭示了签名的空间特异性：
免疫抑制相关基因模块（FOXP3、IL2RA、CTLA4、TGFB1、IL10 等）
在 stroma\_immune 区域和 tumor\_edge 区域表达最高，
而在 tumor\_core 中最低，与"肿瘤-免疫交界面是免疫抑制最活跃位点"的生物学模型一致~\citep{zheng2017landscape}；
CAF/ECM 模块（FAP、ACTA2、COL1A1）在基质区高表达，
与免疫模块的空间协同分布在热图层次聚类中自然分离，
揭示了签名基因集的双模块内在结构。

通路富集分析（Figure 4C）将上述分子特征映射至功能层面。
TGF-$\beta$ 信号通路、IL-10 信号通路和 T 细胞受体激活调控通路的显著富集（FDR $<$ 0.05）
与 Treg 和 M2 型 TAM 介导免疫抑制的核心分子机制直接对应~\citep{sakaguchi2020regulatory,noy2014tumor}；
细胞外基质组织和胶原纤维化通路的富集反映了 CAF 激活后的 ECM 重塑程序~\citep{kalluri2016biology}；
趋化因子信号通路富集（含 CXCL12--CXCR4 和 CCL22--CCR4 相关基因）
机制上对应 TAM/CAF 定向招募 Treg 的关键通路~\citep{domanska2013targeting,curiel2004specific}；
VEGF 信号通路的同期富集则提示免疫抑制生态位中同时存在促血管生成信号，
与抗血管生成和免疫检查点抑制剂联用方案的生物学理论基础相符~\citep{finn2020atezolizumab}。

\subsection{niche gene program 在独立切片 CHC20 中的空间复现}

将 HCC4R 提取的签名基因集应用于独立验证切片 CHC20，
以双模块评分计算免疫-基质组合 niche 评分。
Figure 5A 中 CHC20 切片出现了与 HCC4R 空间分布模式结构相似的局灶性高评分热点，
且热点主要定位于肿瘤-基质交界区域，呈现连续聚集的斑块状分布，
表明 HCC4R 来源的签名基因集在 CHC20 中仍能有效定位具有相似免疫抑制特征的空间区域。

双模块分解（Figure 5A）显示，免疫模块（FOXP3、IL2RA、CTLA4、CD163 等）
与基质/ECM 模块（FAP、ACTA2、COL1A1、SPP1 等）
在 CHC20 中既有重叠（交界区域）又有空间分离（纯基质区 ECM 更高），
两者联合高值区域对应 combined niche score 峰值，
与 HCC4R 中观察的 Treg--TAM--CAF 三角结构在 CHC20 中得到模块化层面的复现~\citep{kalluri2016biology,sakaguchi2020regulatory}。

Figure 5B 小提琴图对比了两切片中 hcc4r\_signature\_score 的全切片分布，
两者均呈右偏分布且高评分尾部比例具有可比性，
Figure 5C/5D 的组合 niche 评分空间分布与箱线图进一步量化了跨样本一致性。
CHC20 的高分区域与其独立进行 Cellular Neighborhoods 聚类所识别的免疫抑制 cluster 在空间上高度重叠，
支持所识别的 niche 是 HCC 的普遍性组织学特征而非单一切片特异现象~\citep{llovet2021hepatocellular}。
值得关注的是，CHC20 为混合型肝细胞-胆管癌，其免疫微环境组成可能与经典 HCC（HCC4R）存在生物学差异，
这一背景下两切片间的评分一致性尤为意义重大。

\subsection{TCGA-LIHC 队列的独立预后验证}

将 HCC4R 空间签名基因集通过 ssGSEA 投影至 TCGA-LIHC 队列370例患者，
Kaplan-Meier 分析（Figure KM）显示高 ssGSEA 评分组（High）的总生存期短于低评分组，
log-rank $p < 0.05$，表明免疫抑制生态位活性越强的患者预后越差，
在流行病学层面支持了空间 niche 具有负性预后价值的生物学假说~\citep{zheng2017landscape,curiel2004specific}。

多变量 Cox 比例风险回归（Figure Cox Forest Plot）在控制年龄和肿瘤分期后，
ssGSEA 评分的风险比 HR $>$ 1（95\% CI 不跨越1），
表明免疫抑制生态位签名是独立于临床分期之外的不良预后因子~\citep{cox1972regression}。
若 HR 效应量偏小（HR $\approx$ 1.3）或在调整分期后 $p$ 值趋于边缘显著，
一个合理的解释是 ssGSEA 评分与肿瘤分期之间存在一定共线性——
免疫抑制信号在晚期肿瘤中自然偏高，多变量模型将分期纳入后会部分"调整"评分的独立贡献~\citep{tomczak2015cancer}。
此外，从 Visium 局部空间信号推导的基因签名在 bulk RNA-seq（全肿瘤均质化）中必然经历信号稀释，
这是所有基于空间特征推导临床生物标志物研究的共同局限，
未来以单细胞水平队列数据进行验证有望进一步提高效应量~\citep{barbie2009systematic}。

% ============================================================
%  第四部分（附录）：代码实现细节
% ============================================================

\newpage
\appendix
\renewcommand{\thesection}{\Alph{section}}

\section{配体-受体通讯对列表}
\label{app:lr_table}

\begin{table}[H]
\centering
\caption{纳入分析的配体-受体通讯对}
\label{tab:lr_pairs}
\begin{tabularx}{\textwidth}{lllX}
\toprule
\textbf{配体} & \textbf{受体} & \textbf{通讯对} & \textbf{生物学意义} \\
\midrule
CXCL12 & CXCR4   & CXCL12--CXCR4  & CAF 分泌，招募 Treg/MDSC~\citep{domanska2013targeting} \\
TGFB1  & TGFBR1  & TGFB1--TGFBR1  & 免疫抑制细胞因子，CAF+Treg 共分泌~\citep{mariathasan2018tgfb} \\
PDCD1  & CD274   & PD-1--PD-L1    & 经典免疫检查点轴~\citep{sharma2023immunosuppressive} \\
TIGIT  & NECTIN2 & TIGIT--NECTIN2 & Treg 抑制性检查点~\citep{sakaguchi2020regulatory} \\
LAG3   & HLA-DRA & LAG3--MHC-II   & T 细胞耗竭相关~\citep{sharma2023immunosuppressive} \\
CCL22  & CCR4    & CCL22--CCR4    & TAM 招募 Treg~\citep{curiel2004specific} \\
IL10   & IL10RA  & IL10--IL10RA   & Treg 免疫抑制细胞因子轴~\citep{noy2014tumor} \\
SPP1   & CD44    & SPP1--CD44     & TAM 分泌骨桥蛋白~\citep{biswas2012macrophage} \\
MIF    & CD74    & MIF--CD74      & 巨噬细胞迁移抑制因子~\citep{noy2014tumor} \\
VEGFA  & KDR     & VEGFA--KDR     & 血管生成，CAF 相关~\citep{kalluri2016biology} \\
LGALS9 & HAVCR2  & Galectin9--TIM-3 & 抑制性检查点~\citep{sharma2023immunosuppressive} \\
CCL2   & CCR2    & CCL2--CCR2     & 巨噬细胞招募通路~\citep{biswas2012macrophage} \\
\bottomrule
\end{tabularx}
\end{table}

\section{辅助评分公式}
\label{app:scores}

辅助评分体系定义如下，其中 $z(\cdot)$ 表示跨 spot Z-score 标准化：

\begin{align}
S_{\text{Treg-like},i} &= z(\hat{p}_{\text{Treg},i}) + z(\hat{s}_{\text{gene},i}) \\
S_{\text{immune-stroma},i} &= z(\hat{p}_{\text{Treg},i}) + z(\hat{p}_{\text{T/NK},i}) + z(\hat{p}_{\text{Myeloid},i}) + z(\hat{p}_{\text{Fibroblast},i}) \\
S_{\text{niche},i} &= z(\hat{p}_{\text{Treg},i}) + z(\hat{p}_{\text{Myeloid},i}) + z(\hat{p}_{\text{Fibroblast},i}) + z(\hat{s}_{\text{gene},i}) + z(\bar{n}_{\text{Hep},i})
\end{align}

其中 $\hat{p}$ 为 Cell2location 后验丰度归一化比例，$\hat{s}_{\text{gene}}$ 为免疫抑制基因模块评分，
$\bar{n}_{\text{Hep}}$ 为 spot 邻域内肝细胞比例均值（用于捕捉肿瘤-基质界面效应）。
双模块评分（用于 CHC20 验证）：

\begin{align}
\text{immune\_score}_i &= \frac{1}{|G_{\text{immune}}|}\sum_{g \in G_{\text{immune}}} E_{ig} \\
\text{stromal\_score}_i &= \frac{1}{|G_{\text{stromal}}|}\sum_{g \in G_{\text{stromal}}} E_{ig} \\
\text{combined\_niche\_score}_i &= z(\text{immune\_score}_i) + z(\text{stromal\_score}_i)
\end{align}

$E_{ig}$ 为 log1p 标准化（CP10K）表达量；免疫模块21个基因涵盖 FOXP3、IL2RA、CTLA4、CD163 等，
基质/ECM 模块21个基因涵盖 FAP、ACTA2、COL1A1、SPP1 等。

\section{核心代码实现}
\label{app:code}

\subsection{Cell2location 反卷积}

\begin{lstlisting}[caption={RegressionModel 训练与空间建模核心代码}]
# Stage 1: 训练参考签名矩阵
from cell2location.models import RegressionModel
RegressionModel.setup_anndata(adata_sc, labels_key='final_celltype')
mod = RegressionModel(adata_sc)
mod.train(max_epochs=250, batch_size=1024, use_gpu=True,
          early_stopping=True, early_stopping_patience=30)
cell_state_df = mod.export_posterior(adata_sc,
    use_cols=mod.feature_in_columns)['means_per_cluster_mu_fg']

# Stage 2: 空间建模
import cell2location
cell2location.models.Cell2location.setup_anndata(adata_vis)
mod = cell2location.models.Cell2location(adata_vis,
    cell_state_df=cell_state_df,
    N_cells_per_location=30, detection_alpha=20)
mod.train(max_epochs=1000, batch_size=1000, use_gpu=True)
adata_vis = mod.export_posterior(adata_vis,
    sample_kwargs={'num_samples': 1000, 'batch_size': 1000})
# 结果存储在 adata_vis.obsm['means_cell_abundance_w_sf']
\end{lstlisting}

\subsection{邻域组成聚类与生态位注释}

\begin{lstlisting}[caption={邻域组成向量计算、Leiden 聚类与综合排名注释}]
from scipy.spatial import KDTree
import scanpy as sc

# 构建 kNN 邻接图并计算邻域组成向量
tree = KDTree(coords)
_, indices = tree.query(coords, k=k + 1)  # 排除自身
neighbors = [indices[i, 1:] for i in range(len(coords))]
arr = proportions.to_numpy()
neighborhood_comp = np.array(
    [arr[idx].mean(axis=0) if len(idx) else arr[i]
     for i, idx in enumerate(neighbors)])

# Leiden 聚类（直接在组成向量特征空间，无需 PCA）
adata_nc = ad.AnnData(neighborhood_comp.astype(float))
sc.pp.neighbors(adata_nc, n_neighbors=n_neighbors, use_rep='X')
sc.tl.leiden(adata_nc, resolution=0.5, key_added='nc_cluster')

# 综合排名法注释免疫抑制 niche
stats = df.groupby('nc_cluster')[
    [treg_col, myeloid_col, fibroblast_col, gene_score_col]].mean()
for col in stats.columns:
    stats[f'rank_{col}'] = stats[col].rank(ascending=False)
stats['combined_rank'] = stats.filter(like='rank_').sum(axis=1)
niche_id = stats['combined_rank'].idxmin()
\end{lstlisting}

\subsection{三层特征基因筛选}

\begin{lstlisting}[caption={Layer 1 差异检验、Layer 3 Gini Index 计算}]
# Layer 1: block-level Mann-Whitney U + AUC
from scipy.stats import mannwhitneyu
# 空间网格聚合（降低空间自相关）
gx = np.floor(df['spatial_x'].values / block_size).astype(int)
gy = np.floor(df['spatial_y'].values / block_size).astype(int)
block_ids = [f'{x}_{y}' for x, y in zip(gx, gy)]
# 计算 AUC = U / (n_high * n_low)
u_stat, pval = mannwhitneyu(high_vals, low_vals, alternative='greater')
auc = u_stat / (len(high_vals) * len(low_vals))

# Layer 3: Gini Index 筛选局灶性稀疏基因
def gini_index(x):
    x = np.sort(x.ravel())
    n = len(x)
    cumx = np.cumsum(x)
    return (n + 1 - 2 * cumx.sum() / cumx[-1]) / n if cumx[-1] > 0 else 0
\end{lstlisting}

\subsection{双模块评分与 ssGSEA}

\begin{lstlisting}[caption={双模块评分计算（Python）与 ssGSEA 评分（R）}]
# Python: 双模块评分
def compute_dual_module_scores(adata, immune_genes, stromal_genes):
    expr = _expression_frame(adata)  # CP10K + log1p
    immune_score = expr[[g for g in immune_genes
                          if g in expr.columns]].mean(axis=1)
    stromal_score = expr[[g for g in stromal_genes
                           if g in expr.columns]].mean(axis=1)
    def _z(s): return (s - s.mean()) / s.std() if s.std() > 0 else s * 0
    return immune_score, stromal_score, _z(immune_score) + _z(stromal_score)
\end{lstlisting}

\begin{lstlisting}[language=R, caption={ssGSEA + 多变量 Cox 回归（R）}]
library(GSVA); library(survival); library(survminer)
gsva_scores <- gsva(expr_matrix, list(sig = signature_genes),
    method='ssgsea', kcdf='Gaussian', abs.ranking=TRUE, min.sz=5)
cox_model <- coxph(
    Surv(OS_time, OS_event) ~ ssgsea_score + age + stage,
    data = clinical_data)
summary(cox_model)
\end{lstlisting}

% ============================================================
% 参考文献
% ============================================================
\bibliographystyle{unsrt}
\bibliography{report_0701}

\end{document}
